\Askhsh{21838}{2}
\begin{ylh}{1}
    \item 10.3. Εμβαδόν βασικών ευθύγραμμων σχημάτων
    \item 10.4. Άλλοι τύποι για το εμβαδόν τριγώνου.
\end{ylh}
\begin{wrapfigure}[7]{r}{4cm} % Ο αριθμός των γραμμών και το πλάτος μπορεί να χρειαστεί προσαρμογή
    \centering
    \begin{tikzpicture}[scale=0.8]
        % Ορισμός κορυφών A, B, Γ
        \tkzDefPoint(0,5){A}
        \tkzDefPoint(-4,0){B}
        \tkzDefPoint(6,0){C} % Χρησιμοποιούμε C για τη Γ
        
        % Ορισμός μέσων E στην AB και Δ στην AΓ
        \tkzDefMidPoint(A,B) \tkzGetPoint{E}
        \tkzDefMidPoint(A,C) \tkzGetPoint{D} % Χρησιμοποιούμε D για τη Δ
        
        % Σχεδίαση του τριγώνου ΑΒΓ
        \tkzDrawPolygon(A,B,C)
        
        % Σχεδίαση των διαμέσων ΒΔ και ΓΕ
        \tkzDrawSegment(B,D)
        \tkzDrawSegment(C,E)
        
        % Σχεδίαση των σημείων
        \tkzDrawPoints(A,B,C,E,D)
        
        % Ετικέτες των σημείων
        \tkzLabelPoint[above](A){A}
        \tkzLabelPoint[below left](B){B}
        \tkzLabelPoint[below right](C){$\Gamma$}
        \tkzLabelPoint[left](E){E}
        \tkzLabelPoint[right](D){$\Delta$}
    \end{tikzpicture}
\end{wrapfigure}
Δίνεται τρίγωνο $\text{AB}\Gamma$ με πλευρές $\text{AB}=8$, $\text{A}\Gamma=12$ και γωνία $\hat{\text{A}} = 60^\circ$.
\begin{alist}
    \item Να αποδείξετε ότι το εμβαδόν του τριγώνου $\text{AB}\Gamma$ είναι $(\text{AB}\Gamma)=24\sqrt{3}$. \monades{13}
    \item Αν $\text{B}\Delta$ και $\Gamma\text{E}$ διάμεσοι του τριγώνου $\text{AB}\Gamma$, να αποδείξετε ότι:
    \begin{rlist}
        \item Τα τρίγωνα $\text{B}\text{E}\Gamma$ και $\text{A}\text{E}\Gamma$ είναι ισοδύναμα. \monades{4}
        \item Τα τρίγωνα $\text{E}\text{B}\Gamma$ και $\Delta\text{B}\Gamma$ είναι ισοδύναμα με $(\text{E}\text{B}\Gamma)=(\Delta\text{B}\Gamma)=12\sqrt{3}$. \monades{8}
    \end{rlist}
\end{alist}